\section{Autonomous Runtime State Machine \& LangGraph Architecture}
\label{sec:runtime_state_graph}

\subsection{The 8-Stage Cyclical Execution Graph}
Traditional linear agent workflows (e.g., standard LangChain chains or single-turn ReAct loops) are fundamentally incapable of robust enterprise execution: they cannot handle state rollback, dynamic re-planning on environmental faults, deterministic policy interrupts, or independent multi-path verification.

CentraOperator implements an **8-stage cyclical state graph** powered by **LangGraph** with dynamic conditional branching and feedback edges:

\begin{figure}[htbp]
  \centering
  \begin{tikzpicture}[
  scale=0.92, transform shape,
  node distance=1.1cm and 1.5cm,
  state/.style={rectangle, rounded corners=2.5mm, draw=centrablue, fill=centrabg, line width=1pt, minimum width=2.4cm, minimum height=0.9cm, align=center, font=\small\sffamily\bfseries},
  startstop/.style={ellipse, draw=centradark, fill=centradark!10, line width=1pt, minimum width=2.2cm, minimum height=0.8cm, align=center, font=\small\sffamily\bfseries},
  decision/.style={diamond, draw=centragold, fill=centragold!10, line width=1pt, inner sep=1.5pt, align=center, font=\scriptsize\sffamily\bfseries, aspect=1.6},
  hitl/.style={rectangle, rounded corners=2.5mm, draw=centrared, fill=centrared!10, line width=1.2pt, minimum width=2.6cm, minimum height=0.9cm, align=center, font=\small\sffamily\bfseries},
  edge/.style={->, >=Stealth, line width=1pt, draw=centrablue},
  rededge/.style={->, >=Stealth, line width=1.2pt, draw=centrared},
  greenedge/.style={->, >=Stealth, line width=1.2pt, draw=centragreen},
  dashededge/.style={->, >=Stealth, line width=1pt, draw=centrateal, dashed}
]

% Top Row: Ingestion to Gating
\node[startstop] (goal) {Goal Ingest};
\node[state, right=1.0cm of goal] (understand) {Understand\\Node};
\node[state, right=1.0cm of understand] (recall) {Recall\\Node};
\node[state, right=1.0cm of recall] (plan) {Plan\\Node};
\node[state, below=1.4cm of plan] (gate) {Gate\\Node};

% HITL Interrupt Branch
\node[hitl, right=1.4cm of gate] (interrupt) {HITL Interrupt\\Pause \& Card};
\node[decision, below=1.1cm of interrupt] (human) {Human\\Approval?};

% Execution and Observation (Middle/Bottom Row)
\node[state, left=1.3cm of gate] (act) {Act\\Node};
\node[state, left=1.3cm of act] (observe) {Observe\\Node};
\node[state, below=1.6cm of act] (adapt) {Adapt\\Node};

% Verification and Completion
\node[state, below=1.6cm of observe] (verify) {Verify\\Node (Read-Path)};
\node[startstop, draw=centragreen, fill=centragreen!15, left=1.2cm of verify] (complete) {Complete\\Verified Pass};

% Normal Forward Flow
\draw[edge] (goal) -- (understand);
\draw[edge] (understand) -- (recall);
\draw[edge] (recall) -- (plan);
\draw[edge] (plan) -- (gate) node[midway, right, font=\tiny] {Action Selected};
\draw[greenedge] (gate) -- (act) node[midway, above, font=\tiny\bfseries] {Policy Cleared};
\draw[edge] (act) -- (observe) node[midway, above, font=\tiny] {DOM Model};

% Decision at Observe
\draw[edge] (observe) |- (plan) node[pos=0.25, left, font=\tiny] {Step Success};
\draw[rededge] (observe) -- (adapt) node[midway, left, font=\tiny\bfseries\color{centrared}] {Fault / Anomaly};

% Gate HITL Pauses
\draw[rededge] (gate) -- (interrupt) node[midway, above, font=\tiny\bfseries\color{centrared}] {Threshold Met};
\draw[rededge] (interrupt) -- (human);
\draw[greenedge] (human) -| (act) node[pos=0.2, above, font=\tiny\bfseries\color{centragreen}] {Approved (Resume)};
\draw[rededge] (human.south) -- ++(0,-0.4) -| (adapt) node[pos=0.15, below, font=\tiny\bfseries\color{centrared}] {Rejected};

% Verification Flow
\draw[edge] (plan) -- ++(0,-3.8) -| (verify) node[pos=0.1, above, font=\tiny] {All Plan Steps Executed};
\draw[greenedge] (verify) -- (complete) node[midway, above, font=\tiny\bfseries\color{centragreen}] {API Matches State};
\draw[rededge] (verify) -- (adapt) node[midway, above, font=\tiny\bfseries\color{centrared}] {Read Failure / Lying UI};

% Dynamic Re-planning Loop
\draw[dashededge] (adapt) -- (plan) node[midway, left, font=\tiny\bfseries\color{centrateal}] {Mutated Strategy / Backoff};

\end{tikzpicture}

  \caption{The CentraOperator 8-Stage Cyclical State Machine. Feedback edges allow dynamic re-planning from both the Observe node (on environmental faults) and the Verify node (on independent read failures).}
  \label{fig:graph_cyclical}
\end{figure}

The execution sequence is defined as follows:
\begin{equation}
  \text{Goal} \xrightarrow{} \text{Understand} \xrightarrow{} \text{Recall} \xrightarrow{} \text{Plan} \xrightarrow{} \text{Gate} \xrightarrow{} \text{Act} \xrightarrow{} \text{Observe} \xrightarrow{} \text{Adapt} \xrightarrow{} \text{Verify} \xrightarrow{} \text{Complete}
\end{equation}

Crucially, the graph features three dynamic backward feedback edges:
\begin{enumerate}
  \item $\text{Observe} \xrightarrow{\text{Fault Detected}} \text{Adapt} \xrightarrow{} \text{Plan}$: When an environmental anomaly (e.g., expired session, interstitial banner) occurs, the plan is mutated on the fly without halting execution.
  \item $\text{Verify} \xrightarrow{\text{Verification Failed}} \text{Adapt} \xrightarrow{} \text{Plan}$: When the target system claims success in the UI but the independent read-path discovers that no record settled in the database, the anomaly is flagged, retried, or halted.
  \item $\text{Gate} \xrightarrow{\text{Threshold Exceeded}} \text{HITL Interrupt} \xrightarrow{\text{Human Action}} \text{Resume}$: When an action violates safety criteria, the graph serializes its state, pauses execution, and yields control to an asynchronous approval queue.
\end{enumerate}

\subsection{The Generic RunState Schema}
A paramount design achievement is that the runtime state contains **zero domain-specific hardcoding**. It handles procurement, ledger aging, vendor management, or general IT provisioning through the generalized \texttt{RunState} schema:

\begin{lstlisting}[language=Python, caption={Generic RunState Schema in Python (\texttt{services/runtime/graph/state.py})}]
from typing import TypedDict, List, Dict, Any, Optional, Annotated
import operator

class RunState(TypedDict):
    task_id: str
    company_id: str
    goal: str
    criteria: List[Dict[str, Any]]             # acceptance checks derived in understand
    facts: Dict[str, Any]                      # extracted fields, IDs, page observations
    policies: List[Dict[str, Any]]             # SOPs retrieved from Qdrant with schemas
    memories: List[Dict[str, Any]]             # past human decisions retrieved from Qdrant
    plan: List[Dict[str, Any]]                 # revisable list of intended operations
    pending_action: Optional[Dict[str, Any]]   # next tool invocation to be gated
    last_observation: Optional[Dict[str, Any]] # output of previous tool execution
    history: Annotated[List[Dict[str, Any]], operator.add] # immutable audit event log
    attempts: Dict[str, int]                   # per-step retry and anomaly counters
    budget: Dict[str, int]                     # maximum steps, retries, and LLM tokens
    verification: Optional[Dict[str, Any]]     # read-path verification outcomes
    status: str                                # current graph state machine status
    evidence: List[Dict[str, Any]]             # verifiable proof objects collected
    interrupt_data: Optional[Dict[str, Any]]   # HITL context payload when paused
\end{lstlisting}

\subsection{Detailed Node Transition Mechanics}
Each node in the state graph is a pure Python transformation function:

\subsubsection{1. Understand Node (\texttt{understand.py})}
Deconstructs the raw user prompt into:
\begin{itemize}
  \item \textbf{Intent Classification:} Determines the operational workflow (e.g., \texttt{invoice\_processing}, \texttt{aging\_review}).
  \item \textbf{Criteria Generation:} Establishes explicit acceptance checks (e.g., ``Invoice record exists in GL'', ``PO line item matched'', ``Status is POSTED'').
  \item \textbf{Entity Extraction:} Extracts initial hints such as vendor names, invoice identifiers, and date boundaries.
\end{itemize}

\subsubsection{2. Recall Node (\texttt{recall.py})}
Queries persistent organizational memory:
\begin{itemize}
  \item \textbf{SOP Retrieval:} Performs vector similarity search against Qdrant to retrieve official corporate SOPs applicable to the classified intent.
  \item \textbf{Episodic Memory Retrieval:} Retrieves past human approvals, rejections, and managerial overrides for the specific vendor or operation.
\end{itemize}

\subsubsection{3. Plan Node (\texttt{plan.py})}
Constructs and maintains an ordered sequence of actions. When invoked after \texttt{Observe} or \texttt{Adapt}, it inspects \texttt{facts} and \texttt{last\_observation} to either progress to the next step, select a tool, or route to \texttt{Verify} once all planned actions have finished.

\subsubsection{4. Gate Node (\texttt{gate.py})}
The deterministic policy enforcement boundary. Discussed in Section~\ref{sec:deterministic_gating}.

\subsubsection{5. Act Node (\texttt{act.py})}
Dispatches the approved tool from the \texttt{ToolRegistry}, handles idempotency locking via Redis SHA-256 keys, injects credentials, and executes the network request.

\subsubsection{6. Observe Node (\texttt{observe.py})}
Ingests the raw tool output (HTML DOM or JSON response), transforms it into a structured semantic model, updates \texttt{facts}, and executes anomaly detection routines (identifying CSRF errors, expired cookies, or interstitial banners).

\subsubsection{7. Adapt Node (\texttt{adapt.py})}
The fault recovery engine. Discussed in Section~\ref{sec:fault_recovery}.

\subsubsection{8. Verify Node (\texttt{verify.py})}
Independent read-path verification engine. Discussed in Section~\ref{sec:independent_verifier}.

\subsubsection{9. Complete Node (\texttt{complete.py})}
Finalizes the task, verifies the cryptographic integrity of the SHA-256 audit chain, formats verified evidence artifacts, and issues a final completion summary.

\subsection{State Checkpointing and Durable Pause/Resume}
When an action requires human approval, the workflow cannot block an operating system thread or hold an open HTTP connection. 

CentraOperator utilizes a durable checkpointing mechanism:
\begin{lstlisting}[language=Python, caption={Non-Blocking State Checkpoint \& Resumption (\texttt{services/runtime/graph/workflow.py})}]
# When HITL interrupt triggers in gate_node:
CHECKPOINTS[state["task_id"]] = {
    "state": state,
    "audit": audit,
    "paused_at": time.time()
}
audit.append_event("graph_interrupt_yielded", {
    "task_id": state["task_id"],
    "reason": state.get("interrupt_data", {}).get("reason")
})
return state  # Returns immediately with status="AWAITING_APPROVAL"

def resume_task(self, task_id: str, human_decision: Dict[str, Any]) -> RunState:
    ckpt = CHECKPOINTS[task_id]
    state: RunState = ckpt["state"]
    audit: AuditLogger = ckpt["audit"]
    state["facts"]["human_decision"] = human_decision
    audit.append_event("human_decision_resumed", human_decision)
    del CHECKPOINTS[task_id]
    return self.execute_turn(state, audit)
\end{lstlisting}
This architecture ensures that tasks can remain suspended for minutes, days, or weeks without consuming runtime compute resources, resuming with complete fidelity upon receiving the human decision.
